\theory{Topos theory}{Can one build a mathematical universe that behaves like sets, but whose notion of space, logic, and truth is adapted to a particular problem?}{A topos is a category behaving like sheaves of sets on a space or site. Grothendieck topoi organize geometry and cohomology; elementary topoi provide an alternative foundation with internal logic that may be constructive rather than classical.}{Algebraic geometry, sheaf cohomology, logic, constructive mathematics, category theory, type theory, and theoretical computer science.}{Objects, morphisms, sheaves, and internal logic provide a generalized setting in which geometry and truth can vary together.}

\paragraph{Origin and intuition}
Topos theory arose when Grothendieck combined sheaves with categorical operations in algebraic geometry. Its important application is etale cohomology, which helped prove the Weil conjectures. Later Lawvere and Tierney emphasized topoi as logical universes; Olivia Caramello developed bridge methods between theories and topoi \cite{topos_ref1}.

The category of sets is the familiar example. A topos has objects, maps, products, function objects, a terminal object, and a subobject classifier. It therefore supports much of ordinary mathematics internally, but its logic can differ from classical set theory.

\end{historylist}
\section*{3. Theory}
\subsection*{Core theory}
\paragraph{Grothendieck and elementary topoi}
The category of sheaves of sets on a topological space is a Grothendieck topos. More generally, sheaves on a site form a topos, even when the covering objects are not ordinary open subsets. This lets algebraic geometers choose the topology best suited to the problem.

An elementary topos is defined by finite limits, exponentials, and a subobject classifier. It can interpret intuitionistic logic, so the law of excluded middle need not hold. A topos of $G$-sets is useful when a group symmetry should be built into every object \cite{topos_ref2}.

\paragraph{Concrete illustration: $G$-sets for a simple group}
Let $G=\mathbb{Z}/2\mathbb{Z}=\{0,1\}$ act on a set $X=\{a,b\}$ by swapping $a$ and $b$: the non-identity element of $G$ sends $a\mapsto b$ and $b\mapsto a$. A $G$-set is any set with such a $G$-action; a morphism of $G$-sets is a map $f\colon X\to Y$ satisfying $f(g\cdot x)=g\cdot f(x)$ for all $g\in G$, $x\in X$.

The category $\mathbf{Set}^{G}$ of all $G$-sets is a topos. Its terminal object is the one-element set with trivial $G$-action. Products are cartesian products with $G$ acting componentwise: $(X\times Y)_g=(x,y)\mapsto(g\cdot x,g\cdot y)$. The subobject classifier is $\Omega=\{0,1\}$ where $G$ acts trivially on both elements: the characteristic map of the $G$-invariant subset $\{a\}\subset\{a,b\}$ sends both $a$ and $b$ to $1$, reflecting that $\{a\}$ is not $G$-stable and so its ``truth value'' is the whole of $\Omega$.

This tiny example shows the topos machinery in action: internal logic (truth values in $\Omega$), subobject classification (invariant subsets), and the failure of excluded middle, since the proposition ``$x=a$'' has no definite truth value when $G$ may swap the points \cite{topos_ref1,topos_ref2}.

\paragraph{Logic and applications}
The internal language of a topos treats objects as types or sets and arrows as functions. A subobject is a predicate, and its classifier behaves like a truth-value object. In a Boolean topos the internal logic is classical; in a general topos it is intuitionistic. This makes topoi useful for constructive mathematics and semantics of computation.

Topoi connect category theory with algebraic geometry, homotopy theory, type theory, and logic. Classifying topoi represent geometric theories, while higher topoi extend the idea to homotopy-coherent spaces \cite{topos_ref3}.

\section*{4. Important theorems and results}
\begin{enumerate}[leftmargin=*,itemsep=0.35em]
\item \textbf{Topos sheaf theorem.} Sheaves on a site form a Grothendieck topos.

\item \textbf{Subobject classifier.} In a topos, subobjects of an object correspond to maps into a distinguished truth-value object.

\item \textbf{Exponential object.} Maps into an exponential $B^A$ correspond naturally to maps from a product with $A$ into $B$.

\item \textbf{Lawvere--Tierney topology.} Local notions of truth and sheafification inside a topos are described by suitable closure operators.

\item \textbf{Giraud principle.} Grothendieck topoi can be characterized abstractly by categorical exactness and generator conditions.

These results explain a topos at three levels. The sheaf theorem supplies the objects, the subobject classifier supplies an internal notion of truth, and exponentials supply an internal function space. Lawvere--Tierney topologies describe controlled changes of truth or sheafification. Giraud's principle is the recognition theorem: it tells when an abstract category has enough structure to behave like a category of sheaves.
\end{enumerate}
\noindent\emph{Source for these results:} \cite{topos_ref1}.

\theoryapplications
\theoryconnections{topos_theory}{%
\item Category theory supplies objects, arrows, and universal constructions, while sheaves attach local data to a space and specify how that data restricts and glues.
\item Internal logic allows the topos itself to determine what counts as a truth value, and algebraic geometry supplies sites whose covers reflect arithmetic rather than ordinary open sets.
\item Type theory gives a computational reading of the same logical structure \cite{topos_ref1,topos_ref2}.
}{%
\item Topos theory helps cohomology because sheaf data can be resolved and its failure to glue can be measured.
\item It helps constructive mathematics by providing a setting where excluded middle need not hold, and helps computer science by interpreting types, programs, and proofs in a common semantics.
\item Homotopy-theoretic topoi extend this idea by treating equivalent proofs or paths as geometric data rather than identifying everything strictly \cite{topos_ref1,topos_ref3}.
\item \textbf{Temporal logic and software verification.} In computer science, a program's execution over discrete time steps can be modeled as a $G$-set where $G=\mathbb{Z}$ acts by shifting one step forward. The topos of $\mathbb{Z}$-sets provides an internal logic in which propositions such as ``eventually $P$ holds'' are interpreted by quantifying over the orbit of a state under the $\mathbb{Z}$-action. Model-checking tools for linear temporal logic (LTL) exploit this connection: the truth value of a temporal formula at a state is determined by the entire forward orbit, which is exactly a $G$-set morphism into the subobject classifier of $\mathbf{Set}^{\mathbb{Z}}$ \cite{topos_ref1,topos_ref3}.
\item \textbf{Type theory and proof assistants.} The Curry--Howard correspondence identifies types with propositions and programs with proofs. An elementary topos provides a model of dependent type theory: function types correspond to exponential objects, product types to categorical products, and the unit type to the terminal object. Proof assistants such as Coq and Agda are based on type theories whose semantics can be interpreted in topoi, ensuring that a verified program is correct not only syntactically but also in a rich mathematical universe \cite{topos_ref2,topos_ref3}.
}
\section*{Glossary}
\begin{description}[leftmargin=*,style=nextline,itemsep=0.3em]
\item[\textbf{Topos.}] A category that behaves like the category of sets, supporting internal logic and geometric constructions.
\item[\textbf{Grothendieck topos.}] A category equivalent to the category of sheaves on a site; it organizes geometry and cohomology.
\item[\textbf{Elementary topos.}] A category with finite limits, exponentials, and a subobject classifier; it provides an alternative foundation for mathematics.
\item[\textbf{Subobject classifier.}] A distinguished object that classifies monomorphisms via characteristic maps, serving as a truth-value object in a topos.
\item[\textbf{Exponential object.}] An object $B^A$ representing functions from $A$ to $B$; it supports internal function spaces.
\item[\textbf{Internal logic.}] The logic interpreted within a topos, which may be intuitionistic rather than classical.
\item[\textbf{Sheaf.}] A contravariant functor from a site to sets that satisfies a gluing condition for local data.
\item[\textbf{Site.}] A category equipped with a Grothendieck topology that specifies covering families for defining sheaves.
\item[\textbf{Terminal object.}] An object $1$ such that for every $X$ there is exactly one morphism $X\to 1$.
\item[\textbf{$G$-set.}] A set equipped with an action by a group $G$; the category of $G$-sets is a topos.
\item[\textbf{Intuitionistic logic.}] A logic in which the law of excluded middle is not assumed; the internal logic of a general topos.
\item[\textbf{Lawvere–Tierney topology.}] A closure operator on subobjects inside a topos defining local truth and generalizing sheafification.
\end{description}
\section*{7. Sample Questions}
\emph{Exercise reference:} \cite{topos_ref1}. The cited edition is the source for the exercise set; consult its Exercises section for the official statement and numbering.
\begin{enumerate}[leftmargin=*,itemsep=0.9em]
\item \textbf{Exercise 1.} In the topos $\mathbf{Set}^G$ of $G$-sets where $G = \mathbb{Z}/2\mathbb{Z}$ acts on $X = \{a, b\}$ by swapping, determine all subobjects of $X$ and their characteristic maps to the subobject classifier $\Omega = \{0, 1\}$ (with trivial $G$-action). \emph{Solution.}$G$-invariant subobjects of $X$ must be $G$-stable subsets. The only $G$-stable subsets are $\varnothing$ and $X = \{a, b\}$. The subset $\{a\}$ is not $G$-stable since $g \cdot a = b \notin \{a\}$. Thus the subobjects are: $\varnothing$ with characteristic map $\chi_\varnothing(a) = \chi_\varnothing(b) = 0$, and $X$ with characteristic map $\chi_X(a) = \chi_X(b) = 1$. This differs from $\mathbf{Set}$, where $X$ would have four subobjects.

\item \textbf{Exercise 2.} In $\mathbf{Set}^G$ with $G = \mathbb{Z}/2\mathbb{Z}$ acting on $X = \{a,b\}$ by swapping, compute the product $X \times X$ as a $G$-set. \emph{Solution.}$X \times X = \{(a,a), (a,b), (b,a), (b,b)\}$ with $G$ acting componentwise: $g \cdot (x,y) = (g \cdot x, g \cdot y)$. Thus $g \cdot (a,a) = (b,b)$, $g \cdot (a,b) = (b,a)$, etc. The orbits are $\{(a,a), (b,b)\}$ and $\{(a,b), (b,a)\}$. The product has two orbits, each of size 2.

\item \textbf{Exercise 3.} In the topos $\mathbf{Set}$, compute the exponential object $2^3$ (where $2 = \{0,1\}$ and $3 = \{0,1,2\}$). \emph{Solution.}$2^3$ is the set of all functions from $\{0,1,2\}$ to $\{0,1\}$. There are $|2|^{|3|} = 2^3 = 8$ such functions. These correspond to the 8 subsets of $\{0,1,2\}$ (via characteristic functions): $\varnothing$, $\{0\}$, $\{1\}$, $\{2\}$, $\{0,1\}$, $\{0,2\}$, $\{1,2\}$, $\{0,1,2\}$.

\item \textbf{Exercise 4.} In the topos of sheaves on a topological space $X$ with the usual open-set topology, explain why the subobject classifier $\Omega$ assigns to each open set $U$ the set of open subsets $V \subseteq U$. \emph{Solution.}$\Omega$ is the sheaf defined by $\Omega(U) = \{V \subseteq U : V \text{ open}\}$. A subobject of a sheaf $\mathcal{F}$ over $U$ corresponds to a collection of sections $\mathcal{F}(V)$ for open $V \subseteq U$, and its characteristic map sends each section $s \in \mathcal{F}(V)$ to the open set $\{x \in V : s(x) \in \text{subobject}_x\}$. In the topological case, truth is local: each point has its own truth value (the open set containing it). This is the key difference from $\mathbf{Set}$, where $\Omega = \{0,1\}$.

\item \textbf{Exercise 5.} Show that in any topos, the product $A \times \Omega$ satisfies $A \times \Omega \cong A$ only if $\Omega \cong 1$ (which fails when the topos is not trivial). \emph{Solution.}$\Omega$ is the subobject classifier. In a topos, $\Omega \cong 1$ if and only if the topos is the degenerate topos (equivalent to $\mathbf{Set}/\{\text{pt}\}$). In any non-degenerate topos, $\Omega$ has at least two distinct global sections (corresponding to $\top$ and $\bot$), so $|\Omega| > 1$ in any model. If $A \times \Omega \cong A$ for all $A$, then taking $A = 1$ gives $\Omega \cong 1$, a contradiction. In non-degenerate topoi, $\Omega \not\cong 1$, and the product $A \times \Omega$ is in general strictly larger than $A$.

\end{enumerate}
\section*{8. References}
\begin{thebibliography}{9}
\bibitem{topos_ref1} S. Mac Lane and I. Moerdijk, \emph{Sheaves in Geometry and Logic}, Springer-Verlag, 1994.
\bibitem{topos_ref2} P. T. Johnstone, \emph{Topos Theory}, Academic Press, 1977.
\bibitem{topos_ref3} O. Caramello, \emph{Theories, Sites, Toposes}, Cambridge University Press, 2018.
\end{thebibliography}
